\documentclass[a4paper,12pt]{article}
\usepackage{amsmath}
\usepackage[russian]{babel}
\usepackage[utf8]{inputenc}
\usepackage[T2A]{fontenc}
\usepackage{geometry}
\usepackage{listings}
\usepackage{xcolor}
\usepackage{graphicx}
\usepackage{color}


\geometry{margin=1.5cm}
% Пример темы в стиле тёмной IDE
\definecolor{intj-bg}{HTML}{2B2B2B}
\definecolor{intj-key}{HTML}{CC7832}
\definecolor{intj-str}{HTML}{6A8759}
\definecolor{intj-comm}{HTML}{808080}
\definecolor{intj-id}{HTML}{A9B7C6}
\definecolor{intj-num}{HTML}{6897BB}

\lstdefinestyle{intellij}{
    language=C++,
    backgroundcolor=\color{intj-bg},
    basicstyle=\ttfamily\small\color{intj-id},
    keywordstyle=\color{intj-key}\bfseries,
    stringstyle=\color{intj-str},
    commentstyle=\itshape\color{intj-comm},
    numberstyle=\color{intj-comm}\tiny,
    stepnumber=1,
    frame=single,
    rulecolor=\color{intj-comm},
    showstringspaces=false,
    breaklines=true,
    tabsize=4
}




\begin{document}

% ---------- Титульный лист ----------
    \begin{titlepage}
        \begin{center}
            \textbf{Федеральное государственное автономное образовательное учреждение высшего образования «Национальный исследовательский университет ИТМО} \\
            \vspace{1cm}

            \textbf{Факультет программной инженерии и компьютерной техники} \\
            \vspace{1cm}

            \Large{\textbf{ОТЧЁТ}} \\
            \Large{\textbf{по задачам: E-H}} \\
            \large{по дисциплине «Алгоритмы и структуры данных»} \\
            \vspace{2cm}
        \end{center}

        \begin{flushright}
            \textbf{Выполнил:} \\
            Ануфриев Андрей Сергеевич \\
            P3219 \\
        \end{flushright}

        \vfill
        \begin{center}
            Санкт-Петербург 2026
        \end{center}
    \end{titlepage}

% ---------- Содержание ----------
    \tableofcontents
    \newpage

% ========== ЗАДАНИЕ 1 ==========


    \section{Задание E. Коровы в стойла}

    \lstinputlisting[style=intellij, label=lst:taskE]{../tasks/E.cpp}

    \subsection*{Суть решения}
    \textbf{Ключевая идея:} использование бинарного поиска по ответу. Ищем нужное расстояние между коровами, которое можно достичь при расстановке k коров в n стойл.\\
    \textbf{Основной алгоритм:}
    \begin{enumerate}
        \item \textbf{Инициализация бин.поиска}: диапазон поиска от 0 до (последнее стойло - первое стойло).
        \item \textbf{Бин.поиск}: для каждого проверяемого расстояния mid проверяем возможность расстановки.
        \item \textbf{Жадно расставляем}:
        \begin{itemize}
            \item первую корову на первое стойло.
            \item Для следующего стойла: если расстоянии >= mid, ставим туда и обновляем last\_pos.
            \item Если удалось расставить, увеличиваем mid.
            \item Если не удалось, уменьшаем.
        \end{itemize}
    \end{enumerate}
    Реализация простого бин.поиска\\
    \textbf{Время: O(n log(max\_dist))}\\
    \textbf{Память: O(n)}


% ========== ЗАДАНИЕ 2 ==========
    \newpage


    \section{Задание F. Число}

    \lstinputlisting[style=intellij, label=lst:taskF]{../tasks/F.cpp}

    \subsection*{Суть решения}
    \textbf{Ключевая идея:} делаем красивую сортировку чисел как букв\\
    \textbf{Почему это работает:} если a+b > b+a, то в любом числе, где a и b стоят рядом, a должна быть перед b для максимизации.\\
    \textbf{Время: O(n log(n L))} - где n - количество фрагментов, L - средняя длина фрагмента.\\
    \textbf{Память: O(S)}

% ========== ЗАДАНИЕ 3 ==========
    \newpage


    \section{Задание G. Кошмар в замке}

    \lstinputlisting[style=intellij,  label=lst:taskG]{../tasks/G.cpp}

    \subsection*{Суть решения}
    \textbf{Ключевая идея:} для каждой буквы с частотой > 1, максимальное расстояние между её вхождениями достигается, когда одно вхождение находится в начале строки, а другое - в конце. Буквы с большим весом должны быть размещены на максимальные расстояния.\\

    \textbf{Основной алгоритм:}
    \begin{enumerate}
        \item Подсчет частот и чтение весов
        \item Создание списка букв
        \item Сортировка по убыванию веса
        \item Используем два указателя: left и right
        \item и расставляем буквы с частотой > 1 по краям, а остальные в середину.
    \end{enumerate}
    Я с начала думал, что учитывается вес каждой буквы(каждого вхождения) и пытался как-то расставлять буквы, которые встречаются >2раз.
    Но потом потыкавшись в тесты, я понял, что надо внимательнее читать задание.\\
    \textbf{Время: O(n + k log k)} - где n - длина строки, k <= 26\\
    \textbf{Память: O(n)}


% ========== ЗАДАНИЕ 4 ==========
    \newpage


    \section{Задание H. Магазин}

    \lstinputlisting[style=intellij, label=lst:taskH]{../tasks/H.cpp}

    \subsection*{Суть решения}
    \textbf{Ключевая идея:} жадный алгоритм с сортировкой.
    \textbf{Основной алгоритм:}
    \begin{enumerate}
        \item Сортировка в порядке убывания
        \item Подсчет суммы с пропусками
    \end{enumerate}
    \textbf{Время: O(n log n)}\\
    \textbf{Память: O(n)}

\end{document}